% 01-2 를 줄이면서 뺀 쪽들 (3주차 정밀도 · 재현율 / ROC · PR-AUC 덱에서 다시 쓸 것)

% --- 회귀 지표
\begin{frame}{회귀의 채점: MAE 와 RMSE}
\vfill
\centering
{\Large
\[ \text{MAE} = \frac{1}{n}\sum_{i=1}^{n} |\,y_i - \hat y_i\,| \qquad\qquad \text{RMSE} = \sqrt{\frac{1}{n}\sum_{i=1}^{n} (y_i - \hat y_i)^2} \]}\vskip0.6em
{\large $y_i$ = 정답, \quad $\hat y_i$ = 예측, \quad 둘 다 \kb{$y$ 와 같은 단위} (m, 개)}\vskip1em
{\large 오차 \texttt{[1, 1, 1, 9]}: \quad MAE $= \dfrac{1+1+1+9}{4} = \kb{3}$ \qquad RMSE $= \sqrt{\dfrac{1+1+1+81}{4}} = \sqrt{21} = \kq{4.58}$}\vskip1em
{\small RMSE 는 제곱하니까 \textbf{크게 틀린 하나}에 더 큰 벌을 준다. 무엇이 더 나쁜지에 따라 고른다}
\vfill
\note{회귀는 얼마나 틀렸는지를 잽니다. MAE는 오차의 절댓값 평균, RMSE는 오차를 제곱해 평균 내고 루트를 씌운 겁니다. 둘 다 y와 단위가 같아서 해석하기 쉽습니다. 오차가 1, 1, 1, 9인 예를 보면 MAE는 3, RMSE는 루트 21, 4.58입니다. 9짜리 오차 하나가 제곱되면서 RMSE를 끌어올렸습니다. 작게 여러 번 틀리는 것보다 한 번 크게 틀리는 게 더 나쁜 문제, 예를 들어 다리 하중 예측이라면 RMSE를 씁니다.}
\end{frame}

% --- 투구 열어 보기
\begin{frame}{투구 데이터 열어 보기: 한 행 = 공 하나}
\vfill
\centering
{\normalsize\renewcommand{\arraystretch}{1.2}
\begin{tabular}{clllll}
\textbf{공} & \textbf{던지기 직전 볼카운트} & \textbf{구종} & \textbf{구속} & \textbf{결과} & \textbf{홈런?}\\\hline
1 & 0-0 & 직구 & 160 km/h & 헛스윙 & 0\\
2 & 0-1 & 체인지업 & 140 km/h & 볼 & 0\\
3 & 1-1 & 체인지업 & 142 km/h & 볼 & 0\\
4 & 2-1 & 직구 & 160 km/h & 헛스윙 & 0\\
5 & 2-2 & 직구 & 161 km/h & 볼 & 0\\
6 & \kq{3-2} & 직구 & 161 km/h & 타격 $\to$ \textbf{홈런} & \kq{1}
\end{tabular}}\vskip0.8em
{\small\color{gray} 2024-06-05 오타니의 한 타석 (Baseball Savant). \textbf{타석} = 타자가 들어서서 결과가 나올 때까지, 공 여러 개 (6월 평균 \hrPerPA 개)}\\[0.6em]
{\normalsize 입력 $X$: \kb{던지기 직전 볼카운트} · 타자 · 투수 \qquad 출력 $y$: 이 공이 홈런인가 (0/1)}\\[0.3em]
{\small\color{gray} 볼카운트는 타석 안에서 공마다 바뀌고, 새 타석이 시작되면 0-0 으로 돌아간다}
\vfill
\note{이제 처음 질문으로 돌아갑니다. 한 행이 공 하나인 데이터를 열어 봅니다. 6월 5일 오타니의 한 타석입니다. 첫 공은 0볼 0스트라이크에서 160킬로 직구, 헛스윙. 두 번째 공은 0볼 1스트라이크에서 체인지업, 볼. 이렇게 볼카운트가 공마다 바뀌다가, 3볼 2스트라이크에서 던진 여섯 번째 공을 쳐서 홈런이 됩니다. 여기서 타석이라는 말을 정리합시다. 타석은 스윙 한 번이 아니라, 타자가 들어서서 아웃이든 안타든 결과가 나올 때까지입니다. 6월에는 한 타석에 공이 평균 \hrPerPA 개였습니다. 한 행이 공 하나이면, 던지기 직전의 볼카운트는 이미 알고 있으니 입력으로 쓸 수 있습니다. 출력 y는 이 공이 홈런이 되었나입니다. 이 표에서는 여섯 행 중 마지막 하나만 1입니다.}
\end{frame}

% --- 볼카운트
\begin{frame}{던지기 직전 볼카운트는 도움이 될까?}
\vfill
\centering\includegraphics[width=0.8\linewidth]{pf_count}\par\vskip0.3em
{\large 0-2 에서는 \kb{0.45\%}, 3-2 에서는 \kq{1.56\%}: 볼카운트에 따라 3배 넘게 다르다}\\[0.3em]
{\small\color{gray} 2024년 6월 MLB 투구 \hrN 개. 3-0 이 낮은 건 타자가 거의 치지 않기 때문 (볼 하나면 볼넷)}
\vfill
\note{던지기 직전에 알 수 있는 볼카운트가 정말 쓸모 있는 입력일까요? 볼카운트별로 홈런이 된 공의 비율을 셌습니다. 투수가 몰린 3볼 2스트라이크, 3볼 1스트라이크에서는 1.5퍼센트가 넘습니다. 투수가 유리한 0볼 2스트라이크에서는 0.45퍼센트입니다. 3배가 넘게 차이 납니다. 투수가 몰리면 스트라이크를 넣어야 하니 치기 좋은 공이 들어오는 겁니다. 3볼 0스트라이크가 가장 낮은 건, 타자가 볼 하나만 더 기다리면 볼넷이라 거의 치지 않기 때문입니다. 볼카운트는 예측하는 순간 알 수 있고, 실제로 차이를 만드는 좋은 입력입니다.}
\end{frame}

% --- 정확도의 함정
\begin{frame}{정확도의 함정: 드문 것을 찾는 문제}
\vfill
\centering
{\large
\begin{tabular}{lcc}
& \textbf{전부 ``홈런 아님''} & \textbf{규칙: 볼카운트가 3-1 · 3-2 면 ``홈런''}\\[0.2em]\hline\\[-0.6em]
정확도 & \kq{\hrAccZero\%} & \paAccRule\%\\[0.5em]
잡아낸 홈런 & \kq{0 개} & \kb{\paTP 개}
\end{tabular}}\vskip1.4em
{\large 홈런을 \kq{하나도 못 잡는} 쪽이 정확도는 더 높다}\vskip0.6em
{\small\color{gray} 홈런이 드물어서, ``아님''을 맞힌 수가 점수를 거의 다 차지한다}
\vfill
\note{처음의 AI와 아주 단순한 규칙 하나를 비교합니다. 왼쪽이 처음 본 AI, 앞에서 본 대로 전부 홈런 아님이라고 답하는 모델입니다. 정확도 \hrAccZero 퍼센트, 잡아낸 홈런은 0개. 오른쪽 규칙은 방금 본 볼카운트를 씁니다. 던지기 직전 볼카운트가 3볼 1스트라이크나 3볼 2스트라이크면 홈런이라고 답합니다. 정확도는 \paAccRule 퍼센트로 오히려 낮은데, 홈런 \paTP 개를 잡아냅니다. 드문 것을 찾는 문제에서 정확도는 거의 아님을 맞힌 수로만 정해지기 때문에, 우리가 원하는 걸 재지 못합니다.}
\end{frame}

% --- 던지기 전에 홈런을 맞힐 수 있나
\begin{frame}{던지기 전 정보로 홈런을 맞힐 수 있을까?}
\markpage{cant}
\vfill
\centering\includegraphics[width=0.7\linewidth]{pf_cant.png}\par\vskip0.3em
{\large 볼카운트 · 타자 · 투수를 다 넣어도 확률은 \kb{두세 배} 오를 뿐, \kq{50\%는 근처에도 못 간다}}\\[0.3em]
{\normalsize 그래서 모델은 모든 공에 ``홈런 아님''. \quad 못 맞히는 문제도 있다 $\to$ \kb{정확도 대신 다른 잣대}로 봐야 한다}\\[0.2em]
{\scriptsize\color{gray} 6/1$\sim$23 학습 · 6/24$\sim$30 시험, 입력 = 볼카운트 · 타자 · 투수 (원-핫), 로지스틱회귀. 홈런 타자 = 6월 1$\sim$23일 공 150개 이상 받은 타자 중 홈런 비율 상위 10\%}
\vfill
\note{학생들이 꼭 묻는 질문입니다. 그럼 볼카운트랑 타자랑 투수를 다 알면 홈런을 맞힐 수 있지 않나요? 실제로 해 봤습니다. 모든 공의 평균 홈런 확률은 0.79퍼센트입니다. 3볼 2스트라이크면 \canTwo 퍼센트, 홈런을 잘 치는 상위 10퍼센트 타자면 \canTop 퍼센트로 두세 배 오릅니다. 셋을 다 넣은 모델이 가장 높게 본 공도 \canMax 퍼센트입니다. 홈런이라고 찍으려면 50퍼센트는 넘어야 하는데 근처에도 못 갑니다. 그래서 모델은 모든 공에 홈런 아님이라고 답하고, 정확도는 99퍼센트가 됩니다. 던지기 전에는 홈런을 맞힐 수 없습니다. 이게 실패가 아니라, 이런 문제가 있다는 걸 알아보는 게 문제 정형화입니다. 맞힐 수 없는 문제에서는 확률이 조금 높은 공을 골라내는 것 자체가 의미가 있고, 그걸 채점하려면 정확도가 아닌 다른 잣대가 필요합니다. 다음 쪽에서 그 잣대를 봅니다. 실제로 잘 맞히는 모델은 다음 시간 선형회귀부터 만듭니다.}
\end{frame}

% --- 세는 법을 바꾼다
\begin{frame}{그래서 세는 법을 바꾼다: 정밀도와 재현율}
\vfill
\centering
{\large\renewcommand{\arraystretch}{1.3}
\begin{tabular}{l|cc}
& \textbf{실제 홈런} & \textbf{실제 아님}\\\hline
\textbf{``홈런''이라고 함} & \kb{\paTP} (잡았다) & \kq{\paFP} (헛짚었다)\\
\textbf{``아님''이라고 함} & \kq{\paFN} (놓쳤다) & (나머지)
\end{tabular}}\vskip1.2em
{\large 재현율 (recall) $= \dfrac{\text{잡았다}}{\text{실제 홈런 전부}} = \dfrac{\paTP}{\paTP + \paFN} = \kb{\paRec\%}$}\vskip0.6em
{\large 정밀도 (precision) $= \dfrac{\text{잡았다}}{\text{``홈런''이라고 한 것 전부}} = \dfrac{\paTP}{\paTP + \paFP} = \kq{\paPrec\%}$}\vskip0.8em
{\small\color{gray} 전부 ``홈런 아님''은 재현율 0\%. 자세한 계산과 둘 사이의 줄다리기는 3주차 (정밀도와 재현율)}
\vfill
\note{그래서 세는 법을 바꿉니다. 실제 홈런 중에서 몇 개를 잡았나, 이게 재현율입니다. 규칙은 \paRec 퍼센트를 잡습니다. 홈런이라고 찍은 것 중에 진짜 홈런이 몇 개였나, 이게 정밀도입니다. \paPrec 퍼센트. 3볼 카운트의 공 \paCalled 개를 전부 홈런이라고 찍었는데, 진짜는 \paTP 개뿐입니다. 치기 좋은 카운트라도 대부분의 공은 홈런이 되지 않으니까요. 전부 아님이라고 하는 모델은 재현율이 0입니다. 정확도 \hrAccZero 퍼센트에 가려 있던 사실이 드러납니다. 이 두 지표는 3주차에 자세히 다룹니다. 오늘은 무엇을 세느냐가 점수의 의미를 바꾼다는 것만 기억하세요.}
\end{frame}

% --- 실습
\begin{frame}[fragile]{실습: 문제를 직접 정해 보기}
\vfill
\centering\begin{tabular}{rl}
\large 1 & 6월 투구: 스윙할까 (분류) · 구속 (회귀), 정확도 · MAE · 기준선 직접 구현\\[0.4em]
\large 2 & 6월 투구 \hrN 개: 전부 ``아님'' \hrAccZero\%, 볼카운트 규칙의 정밀도 · 재현율,\\
& 한 행 = 공 $\to$ 타석 $\to$ 타자 (\hrRate\% $\to$ \hrPArate\% $\to$ \hrBatRate\%)\\[0.4em]
\end{tabular}\vskip0.8em
투구 데이터는 Baseball Savant 에서 받는다 (로그인 없이, 처음엔 5분쯤) \\[0.3em]\href{https://colab.research.google.com/github/SuminHan/book-ml/blob/main/slides/ml1/week01/2_problem_formulation/lab/problem_formulation.ipynb}{\textcolor{ksagolddk}{▶ Colab에서 실습 노트북 열기}}
\vfill
\note{실습에서는 야구 데이터로 직접 다섯 칸을 채웁니다. 같은 투구로 스윙할까라는 분류와 구속이라는 회귀 두 문제를 만들고, 정확도, MAE, 기준선을 직접 함수로 만듭니다. 6월의 모든 투구로 정확도의 함정과 정밀도, 재현율을 계산하고, 한 행을 공에서 타석, 타자로 바꿔 봅니다.}
\end{frame}
